\documentclass[11pt,a4paper]{article}

\usepackage[a4paper,margin=2.4cm]{geometry}
\usepackage[T1]{fontenc}
\usepackage[utf8]{inputenc}
\usepackage{lmodern}
\usepackage{microtype}
\usepackage{xcolor}
\usepackage{booktabs}
\usepackage{longtable}
\usepackage{array}
\usepackage{enumitem}
\usepackage{listings}
\usepackage{hyperref}
\usepackage{fancyhdr}
\usepackage{titlesec}

\definecolor{mohidblue}{RGB}{20,68,114}
\definecolor{mohidgrey}{RGB}{90,90,90}
\definecolor{codebg}{RGB}{247,247,247}

\hypersetup{
  colorlinks=true,
  linkcolor=mohidblue,
  urlcolor=mohidblue,
  citecolor=mohidblue,
  pdftitle={Mohid-NG Coding and Implementation Standard},
  pdfauthor={Mohid-NG Project},
  pdfsubject={Coding standard and implementation discipline for Mohid-NG},
  pdfkeywords={Mohid-NG, coding standard, C++, Voronoi, finite volume, HPC}
}

\setlength{\headheight}{14pt}
\pagestyle{fancy}
\fancyhf{}
\lhead{Mohid-NG Coding and Implementation Standard}
\rhead{Draft}
\cfoot{\thepage}

\titleformat{\section}{\Large\bfseries\color{mohidblue}}{\thesection}{0.7em}{}
\titleformat{\subsection}{\large\bfseries\color{mohidblue}}{\thesubsection}{0.7em}{}
\titleformat{\subsubsection}{\normalsize\bfseries\color{mohidblue}}{\thesubsubsection}{0.7em}{}

\lstdefinestyle{mohidcpp}{
  language=C++,
  basicstyle=\ttfamily\small,
  backgroundcolor=\color{codebg},
  frame=single,
  rulecolor=\color{mohidgrey},
  breaklines=true,
  columns=fullflexible,
  keepspaces=true,
  showstringspaces=false
}

\setlist[itemize]{topsep=3pt,itemsep=2pt,parsep=0pt,leftmargin=1.5em}
\setlist[enumerate]{topsep=3pt,itemsep=2pt,parsep=0pt,leftmargin=1.8em}

\newcommand{\code}[1]{\texttt{#1}}
\newcommand{\shall}{\textbf{shall}}
\newcommand{\mustnot}{\textbf{shall not}}
\newcommand{\should}{\textbf{should}}
\newcommand{\may}{\textbf{may}}

\begin{document}
\pagenumbering{gobble}

\begin{titlepage}
  \centering
  \vspace*{2.0cm}
  {\Huge\bfseries\color{mohidblue} Mohid-NG Coding and Implementation Standard\par}
  \vspace{0.7cm}
  {\Large Formal coding rules for a Voronoi-native environmental modelling platform\par}
  \vspace{1.5cm}
  {\large Draft version for architectural consolidation\par}
  \vfill
  \begin{minipage}{0.86\textwidth}
  \small
  This document defines the coding, organisation, implementation and review standard for Mohid-NG. It is written as a project standard, not as a tutorial. The aim is to support a high-performance, maintainable, scientifically auditable code base without exposing unnecessary implementation complexity to users.
  \end{minipage}
  \vfill
  {\large Mohid-NG Project\par}
\end{titlepage}

\pagenumbering{roman}
\tableofcontents
\newpage
\pagenumbering{arabic}

\section{Purpose and scope}

This document defines how Mohid-NG source code shall be written, organised, reviewed, tested and maintained. It applies to the Mohid-NG solver, numerical modules, physical-model modules, field infrastructure, forcing infrastructure, input and output layers, validation code, examples and public user-facing APIs.

The standard is intended to support the following project goals:

\begin{itemize}
  \item a Voronoi-native finite-volume modelling platform;
  \item high runtime performance on serial and parallel systems;
  \item formal separation between physics, numerics, mesh services, solvers, input and output;
  \item long-term maintainability by a small technical team;
  \item scientific auditability of models, parameters, fields and conservation properties;
  \item a simple public API for users who are not C++ specialists;
  \item advanced C++ implementation internally, where this improves runtime behaviour, safety or extensibility.
\end{itemize}

This standard does not define the mesh-generation algorithms themselves. All mesh generation, remeshing, generator-point movement, Voronoi connectivity construction, mesh-quality analysis, geometry adaptation and mesh-specific partitioning operations belong to VoronoiMeshMaker. Mohid-NG shall consume Voronoi-based computational meshes and mesh metadata through a documented contract.

\section{Architectural principles}

Mohid-NG code shall follow the principles below.

\begin{enumerate}
  \item Public interfaces shall be stable, minimal and physically meaningful.
  \item Implementation details shall remain private to the module that owns them.
  \item Physical models shall not depend on external solver-library types, GIS-library handles, raw file-format handles or VoronoiMeshMaker implementation types.
  \item Mesh-related algorithms shall be implemented in VoronoiMeshMaker, not in Mohid-NG.
  \item The numerical core shall prioritise runtime performance over compilation speed.
  \item Hot loops shall avoid inheritance, virtual functions, heap allocation, string lookup and indirect object graphs.
  \item Data-oriented design shall be the default for topology views, geometry views, fields, particles, face loops, cell loops, layer loops and assembly kernels.
  \item Factories, registries, configuration parsing and input/output code may use higher-level abstractions, provided they do not enter performance-critical loops.
  \item Every physical model shall document its equations, assumptions, parameters, units, fields, diagnostics and validation cases.
  \item Every operation that can alter conserved quantities shall expose diagnostics for conservation error.
\end{enumerate}

The guiding distinction is:

\begin{quote}
Advanced C++ is allowed inside Mohid-NG; advanced C++ shall not be required from the ordinary Mohid-NG user.
\end{quote}

\section{Project boundaries}

\subsection{Mohid-NG}

Mohid-NG owns the physical and numerical modelling platform. It shall provide:

\begin{itemize}
  \item field storage and field registries;
  \item finite-volume operators over Voronoi mesh views;
  \item physical models;
  \item forcing and boundary-condition handling;
  \item time integration;
  \item residual assembly;
  \item solver-backend adapters;
  \item conservation diagnostics;
  \item restart and checkpoint logic;
  \item scientific output and validation infrastructure;
  \item user-facing case execution.
\end{itemize}

Mohid-NG shall not construct or modify Voronoi meshes internally. When mesh adaptation, remeshing or repartitioning is required, Mohid-NG shall request these operations through a formal service contract and shall then receive a new mesh state, mapping metadata, partition metadata and transfer metadata.

\subsection{VoronoiMeshMaker}

VoronoiMeshMaker owns all mesh-specific implementation. This includes:

\begin{itemize}
  \item Voronoi mesh generation in two and three dimensions;
  \item vertically structured Voronoi-column mesh construction, when used;
  \item generator-point insertion, removal and movement;
  \item remeshing;
  \item topology reconstruction;
  \item cell-face-node connectivity;
  \item boundary-patch construction from geometric and GIS inputs;
  \item mesh-quality metrics;
  \item mesh validation;
  \item geometric search structures required by remapping;
  \item mesh hierarchy construction for geometric multigrid, when available;
  \item mesh partitioning metadata, when it is mesh-topology dependent.
\end{itemize}

The interface between Mohid-NG and VoronoiMeshMaker shall be file-based, memory-based or service-based, but it shall remain explicit and documented. Mohid-NG shall not include private VoronoiMeshMaker headers in its numerical kernels.

\section{Public and private structure}

\subsection{Directory organisation}

The recommended organisation is:

\begin{lstlisting}[style=mohidcpp]
include/MohidNG/              // public API
src/MohidNG/                  // implementation
src/MohidNG/Internal/         // private implementation details
tests/                        // unit, integration and validation tests
examples/                     // documented examples
benchmarks/                   // performance and scaling cases
docs/                         // user and developer documentation
\end{lstlisting}

Public headers shall be located under \code{include/MohidNG}. Private implementation headers shall remain under \code{src/MohidNG/Internal} or an equivalent internal path. Public headers shall not include private headers.

\subsection{Public API}

Public API types shall represent stable project concepts. Examples include:

\begin{lstlisting}[style=mohidcpp]
MohidNG::Case
MohidNG::Simulation
MohidNG::MeshView
MohidNG::GeometryState
MohidNG::FieldSet
MohidNG::BoundaryPatch
MohidNG::ForcingSet
MohidNG::SolverOptions
MohidNG::Diagnostics
\end{lstlisting}

Public headers shall not expose:

\begin{itemize}
  \item external numerical backend objects;
  \item raw HDF5, NetCDF, GDAL, OGR or PROJ handles;
  \item private VoronoiMeshMaker implementation types;
  \item internal storage containers not intended as stable API;
  \item implementation-specific work arrays;
  \item raw pointers whose ownership semantics are unclear.
\end{itemize}

\subsection{Internal implementation}

Internal implementation types may use advanced C++ constructs, including templates, concepts, traits, policy classes, static dispatch and compile-time configuration. This complexity shall remain behind stable public APIs or behind clearly documented developer-level extension points.

\section{Naming standard}

\subsection{Namespace}

The project shall use one namespace consistently. The recommended namespace is:

\begin{lstlisting}[style=mohidcpp]
namespace mohidng
\end{lstlisting}

If a capitalised namespace is selected instead, it shall be used consistently throughout the project. Mixing \code{mohidng}, \code{MohidNG}, \code{MohidNg} and other variants is forbidden.

\subsection{Types}

Concrete types shall use \code{PascalCase}:

\begin{lstlisting}[style=mohidcpp]
VoronoiMeshView
GeometryState
CellField
FaceField
FieldSet
BoundaryPatch
TransportModel
SolverOptions
RestartWriter
\end{lstlisting}

Template parameters shall use short but meaningful names. Avoid single-letter template parameters except for conventional mathematical or iterator-like cases.

\subsection{Functions}

Use one function naming convention throughout the project. The recommended convention is \code{snake\_case}.

\begin{lstlisting}[style=mohidcpp]
load_case(...)
load_mesh(...)
advance_time_step(...)
compute_gradient(...)
assemble_residual(...)
write_checkpoint(...)
\end{lstlisting}

Member functions and free functions shall follow the same convention unless there is a documented reason to distinguish them.

\subsection{Enums and constants}

Use scoped enums for typed choices:

\begin{lstlisting}[style=mohidcpp]
enum class FieldLocation {
    Cell,
    Face,
    Node,
    Layer,
    Particle
};
\end{lstlisting}

Avoid unscoped enums in public APIs. Avoid macro constants unless needed for build configuration or third-party integration.

\subsection{Implementation-specific names}

Implementation-specific names shall be explicit when ambiguity is possible:

\begin{lstlisting}[style=mohidcpp]
compute_gradient_wlsqr(...)
assemble_transport_native(...)
write_unstructured_vtu(...)
\end{lstlisting}

This rule shall be applied with restraint. Prefer namespaces, file organisation and type names over excessive function-name length.

\section{Header policy and include discipline}

\subsection{Header guards}

Every header shall use:

\begin{lstlisting}[style=mohidcpp]
#pragma once
\end{lstlisting}

as the first non-comment line.

\subsection{Header contents}

A public header shall contain:

\begin{itemize}
  \item a brief description of its purpose;
  \item the public concepts exposed by the file;
  \item minimal includes;
  \item forward declarations where appropriate;
  \item no private implementation leakage.
\end{itemize}

\subsection{Include order}

Includes shall be ordered as:

\begin{lstlisting}[style=mohidcpp]
#include <standard_library_headers>

#include <external_library_headers>

#include <MohidNG/...>
\end{lstlisting}

Each file shall include only what it uses. Public headers shall avoid heavy external dependencies. External library headers shall be isolated inside implementation files or adapter modules whenever possible.

\section{Formatting}

Mohid-NG shall use automated formatting. A repository-level \code{.clang-format} file shall define the formatting policy.

Rules:

\begin{enumerate}
  \item Every submitted file shall be formatted automatically.
  \item Manual alignment that fights the formatter shall be avoided.
  \item \code{clang-format off} shall be used only in exceptional cases.
  \item Formatting violations shall block merging.
  \item The formatting standard shall apply to production code, tests and examples.
\end{enumerate}

The project should also use \code{clang-tidy} for static checks and may use \code{cmake-format} for CMake files.

\section{Error handling}

\subsection{General rule}

All errors that can occur in production runs shall produce explicit diagnostic information. Silent failure is forbidden.

Errors shall identify:

\begin{itemize}
  \item the operation that failed;
  \item the relevant file, field, mesh entity, model or backend;
  \item the invalid value, if available;
  \item the expected condition;
  \item whether the error is recoverable.
\end{itemize}

\subsection{Exceptions}

Exceptions may be used in non-hot code paths, including:

\begin{itemize}
  \item case parsing;
  \item mesh loading;
  \item input and output;
  \item configuration validation;
  \item model registration;
  \item checkpoint reading;
  \item pre-run validation.
\end{itemize}

Example:

\begin{lstlisting}[style=mohidcpp]
throw MohidError("Mesh file does not contain boundary-patch metadata.");
\end{lstlisting}

\subsection{Hot loops}

Performance-critical kernels shall not throw exceptions inside loops over cells, faces, layers or particles. Validation shall occur before entering such loops. Kernel assumptions shall be documented.

\subsection{Backend errors}

External solver, input/output and GIS-library errors shall be translated into Mohid-NG diagnostics at adapter boundaries. Physical models shall not expose backend-specific error mechanisms.

\section{Argument validation}

Public API functions shall validate their arguments. Validation is mandatory for:

\begin{itemize}
  \item case loading;
  \item mesh loading;
  \item model registration;
  \item field creation;
  \item boundary-condition creation;
  \item forcing interpolation;
  \item restart reading;
  \item solver configuration;
  \item remeshing requests;
  \item field-transfer requests.
\end{itemize}

Validation shall check, where applicable:

\begin{itemize}
  \item file existence and readability;
  \item schema validity;
  \item mesh dimensionality;
  \item Voronoi metadata completeness;
  \item field location and field size;
  \item unit compatibility;
  \item boundary-patch consistency;
  \item time-axis consistency;
  \item coordinate-reference metadata;
  \item parallel ownership and ghost consistency;
  \item solver option consistency.
\end{itemize}

Internal kernels may assume validated input, provided this assumption is stated in the kernel contract.

\section{Data ownership and lifetime}

Mohid-NG shall distinguish between:

\begin{itemize}
  \item owning containers;
  \item non-owning views;
  \item temporary workspaces;
  \item backend-owned data;
  \item external-library handles;
  \item persistent simulation state.
\end{itemize}

Rules:

\begin{enumerate}
  \item Mesh topology and geometry used by physical models shall be exposed through non-owning views.
  \item The simulation state shall own fields, model states and diagnostics.
  \item Solver backend objects shall be owned only by backend adapters.
  \item Raw pointers shall not appear in public APIs unless there is a strong interoperability reason.
  \item Hot loops shall use spans, views or contiguous buffers after ownership has been resolved outside the loop.
  \item A view shall never outlive the object from which it was created.
\end{enumerate}

\section{Data-oriented numerical core}

The numerical core shall be designed for efficient traversal of cells, faces, layers and particles.

Preferred storage patterns:

\begin{itemize}
  \item structure of arrays;
  \item contiguous field values;
  \item integer indices;
  \item offset arrays;
  \item compact adjacency arrays;
  \item face-based loops for fluxes;
  \item cell-based loops for source terms;
  \item layer-based loops for vertical processes;
  \item particle arrays for Lagrangian models.
\end{itemize}

Avoid in hot paths:

\begin{itemize}
  \item one heap-allocated object per cell;
  \item one heap-allocated object per face;
  \item pointer-linked graph structures;
  \item virtual dispatch per entity;
  \item string lookup inside loops;
  \item heap allocation inside loops;
  \item repeated unit conversion inside loops;
  \item repeated geometry lookup through slow maps.
\end{itemize}

\section{Mesh contract}

Mohid-NG shall operate only on Voronoi-based meshes. Mesh construction and mesh alteration are external to Mohid-NG and belong to VoronoiMeshMaker.

\subsection{Accepted mesh information}

A Mohid-NG mesh view shall provide, at minimum:

\begin{itemize}
  \item cell identifiers;
  \item face identifiers;
  \item node identifiers, if required;
  \item cell-face connectivity;
  \item face-cell connectivity;
  \item neighbour lists;
  \item cell volumes or areas;
  \item face areas or lengths;
  \item cell centres;
  \item face centres;
  \item outward face normals;
  \item boundary-patch membership;
  \item region membership;
  \item coordinate-reference metadata;
  \item mesh-quality diagnostics;
  \item partition and ghost metadata, when running in parallel.
\end{itemize}

\subsection{Forbidden mesh implementation inside Mohid-NG}

Mohid-NG shall not implement:

\begin{itemize}
  \item Voronoi tessellation construction;
  \item generator-point insertion;
  \item generator-point removal;
  \item generator-point movement;
  \item remeshing algorithms;
  \item connectivity reconstruction after remeshing;
  \item mesh-quality improvement algorithms;
  \item geometric search structures for mesh generation;
  \item GIS-to-mesh geometric conversion;
  \item mesh coarsening or refinement algorithms;
  \item mesh hierarchy construction.
\end{itemize}

Mohid-NG may request such operations and may provide physical indicators, workload weights or field information needed by VoronoiMeshMaker. The operation itself remains a VoronoiMeshMaker responsibility.

\section{Fields and units}

Every field shall declare:

\begin{itemize}
  \item name;
  \item physical meaning;
  \item location, for example cell, face, node, layer or particle;
  \item unit;
  \item time level;
  \item ownership and ghost status in parallel runs;
  \item whether it is prognostic, diagnostic or temporary;
  \item whether it represents a conserved density, extensive quantity or auxiliary variable.
\end{itemize}

Units shall be mandatory in configuration, input and output metadata. Unit checking should occur outside hot loops. Kernels shall operate on already normalised values in the expected internal unit system.

\section{Physical-model implementation}

New physical models shall be implemented as concrete types satisfying documented model concepts. They shall not be implemented as subclasses of an abstract base class.

Each model shall provide:

\begin{itemize}
  \item model name;
  \item model category;
  \item version;
  \item maintainer;
  \item governing equations;
  \item assumptions and limitations;
  \item input fields;
  \item output fields;
  \item parameters and units;
  \item local source terms;
  \item conservative flux requirements;
  \item boundary-condition requirements;
  \item diagnostics;
  \item validation cases;
  \item minimal examples.
\end{itemize}

Physical models shall not directly modify mesh topology or mesh geometry. If a model requires adaptation, it shall produce indicators or requests that are passed to the mesh-service layer.

\section{Numerical operators}

Numerical operators shall be implemented as reusable services, not hidden inside individual physical models. This includes:

\begin{itemize}
  \item gradient reconstruction;
  \item divergence operators;
  \item face interpolation;
  \item limiters;
  \item diffusive fluxes;
  \item advective fluxes;
  \item source-term projection;
  \item conservative remapping support;
  \item residual assembly support.
\end{itemize}

Gradient reconstruction shall be a first-class numerical service because it is used by physical discretisation, face reconstruction, mesh-adaptation indicators and error diagnostics.

Operator implementations shall document:

\begin{itemize}
  \item required mesh information;
  \item field locations;
  \item expected order of accuracy;
  \item known limitations;
  \item conservation properties;
  \item parallel ghost requirements;
  \item validation tests.
\end{itemize}

\section{Solver backend isolation}

Mohid-NG owns the equations, field semantics, residuals and diagnostics. External numerical backends may solve algebraic systems, nonlinear systems, time-integration subproblems or optimisation subproblems.

Rules:

\begin{enumerate}
  \item Physical models shall not call external solver libraries directly.
  \item Numerical assembly shall produce backend-neutral system views or backend-neutral assembly requests.
  \item Backend adapters shall translate Mohid-NG data into backend-specific data structures.
  \item Backend-specific objects shall not appear in physical-model APIs.
  \item Solver diagnostics shall be returned in backend-neutral form.
\end{enumerate}

Solver diagnostics shall include:

\begin{itemize}
  \item setup time;
  \item solve time;
  \item number of iterations;
  \item residual reduction;
  \item preconditioner rebuild frequency;
  \item memory use;
  \item convergence failures;
  \item parallel efficiency indicators, when available.
\end{itemize}

\section{GIS, coordinate systems and input data}

GIS compatibility is required for realistic environmental simulations, but GIS objects shall not become the internal finite-volume representation.

Rules:

\begin{enumerate}
  \item Coordinate-reference metadata shall be mandatory for georeferenced meshes.
  \item The computational kernel shall operate in a projected metric coordinate system.
  \item GIS handles shall remain inside GIS adapter modules.
  \item Boundary names, region names, forcing zones and monitoring points shall be propagated from GIS inputs into Mohid-NG metadata.
  \item Three-dimensional GIS or surface data shall be converted into VoronoiMeshMaker mesh services before numerical simulation.
\end{enumerate}

Mohid-NG shall store enough metadata to make outputs interpretable in GIS and scientific visualisation tools.

\section{Parallel execution}

Parallel code shall be designed around explicit ownership, ghost data and communication schedules.

Every parallel field shall define:

\begin{itemize}
  \item owned range;
  \item ghost range;
  \item update schedule;
  \item communication pattern;
  \item consistency requirements before each kernel.
\end{itemize}

Dynamic workload changes shall be handled through formal repartitioning and migration mechanisms. Mohid-NG may compute physical and computational workload weights, but mesh-topology-dependent partitioning operations shall be delegated through the mesh-service contract when they require mesh modification or mesh repartition metadata.

\section{Random numbers and reproducibility}

System random-number functions with implementation-dependent behaviour shall not be used in production code. Stochastic simulations shall use approved reproducible generators with explicit seed management.

This applies to:

\begin{itemize}
  \item particle diffusion;
  \item stochastic forcing;
  \item randomised tests;
  \item Monte Carlo experiments;
  \item stochastic parameter sampling.
\end{itemize}

Every stochastic simulation shall write its seed, generator type and relevant distribution parameters to metadata.

\section{Documentation and comments}

Comments shall explain intent, assumptions and invariants. They shall not merely restate obvious code.

Comments should explain:

\begin{itemize}
  \item physical assumptions;
  \item discrete conservation properties;
  \item boundary-condition interpretation;
  \item parallel ownership rules;
  \item remeshing assumptions;
  \item field-transfer guarantees;
  \item solver-preconditioner assumptions;
  \item unit conventions;
  \item coordinate conventions.
\end{itemize}

Comments shall avoid:

\begin{itemize}
  \item dead code;
  \item commented-out blocks with no explanation;
  \item unexplained formulae;
  \item historical notes with no operational value;
  \item debug messages left in production code;
  \item vague claims such as ``optimised'', ``robust'' or ``accurate'' without context.
\end{itemize}

Public documentation shall be written for users. Developer documentation shall be written for maintainers. Scientific documentation shall connect code modules to equations, assumptions and validation cases.

\section{Build configuration and dependencies}

Optional dependencies shall be isolated behind build options and adapter modules. The core shall have a minimal build configuration suitable for development and testing.

Rules:

\begin{enumerate}
  \item Do not spread dependency-specific preprocessor blocks through physical-model code.
  \item Prefer target-level configuration in the build system.
  \item Place backend-specific code in backend-specific translation units.
  \item Document every optional dependency and the feature it enables.
  \item Do not make a heavy optional dependency mandatory unless it is essential to a declared release target.
\end{enumerate}

Use preprocessor conditionals only for optional dependencies, platforms, compiler workarounds and feature availability. Mark the closing condition of long blocks.

\section{Testing standard}

Every new module shall include appropriate tests. The required test type depends on the module.

\begin{longtable}{p{0.28\textwidth}p{0.62\textwidth}}
\toprule
\textbf{Module type} & \textbf{Required tests} \\
\midrule
Field infrastructure & unit tests, size tests, location tests, ghost-update tests \\
Mesh view consumption & consistency tests, boundary-patch tests, geometry-metadata tests \\
Numerical operator & manufactured-solution tests, conservation tests, convergence tests \\
Physical model & equation test, parameter test, validation case, regression case \\
Solver adapter & assembly test, convergence test, failure-path test, diagnostic test \\
Forcing module & interpolation test, unit test, missing-value test, time-axis test \\
Restart module & write-read consistency test, metadata test, restart continuation test \\
Parallel module & ownership test, ghost test, communication test, small MPI regression test \\
Remeshing interface & request validation test, field-transfer diagnostic test, conservation audit \\
\bottomrule
\end{longtable}

No new physical model shall be accepted as official without:

\begin{itemize}
  \item equation documentation;
  \item parameter documentation;
  \item unit declaration;
  \item at least one example case;
  \item validation or regression evidence;
  \item diagnostics;
  \item stated limitations.
\end{itemize}

\section{Examples}

Examples are part of the code standard. They shall be small, reproducible and documented.

Each example shall state:

\begin{itemize}
  \item purpose;
  \item required input files;
  \item expected output files;
  \item physical interpretation;
  \item numerical options used;
  \item how success is verified.
\end{itemize}

Examples shall not be used as dumping grounds for experimental code. Experimental examples shall be clearly marked and shall not be used as validation evidence.

\section{Code review and acceptance}

A change may be accepted only if it satisfies the relevant criteria below:

\begin{itemize}
  \item it builds in the required configurations;
  \item it is formatted automatically;
  \item it passes existing tests;
  \item it adds tests for new behaviour;
  \item it does not expose private implementation details through public headers;
  \item it does not introduce mesh-generation logic into Mohid-NG;
  \item it does not introduce backend-specific types into physical models;
  \item it documents new public APIs;
  \item it documents physical assumptions for new models;
  \item it reports conservation diagnostics when conserved quantities are affected;
  \item it includes an example when the feature is user-facing.
\end{itemize}

\section{Deprecation policy}

Public APIs shall not be removed without a deprecation path unless there is a severe correctness, security or scientific-integrity reason.

Deprecation notices shall identify:

\begin{itemize}
  \item the deprecated API;
  \item the replacement;
  \item the reason for deprecation;
  \item the planned removal stage;
  \item any scientific effect on old case files or outputs.
\end{itemize}

Case-file schemas shall also follow a deprecation policy. Old cases shall either remain readable or fail with an actionable migration message.

\section{Definition of done}

A feature is not complete when it merely compiles. A Mohid-NG feature is complete when:

\begin{enumerate}
  \item the public interface is defined;
  \item the implementation is isolated in the correct module;
  \item ownership and lifetime are clear;
  \item relevant arguments are validated;
  \item errors produce useful diagnostics;
  \item tests cover normal and failure paths;
  \item examples exist when the feature is user-facing;
  \item documentation describes use, assumptions and limitations;
  \item conservation diagnostics exist when conserved quantities are affected;
  \item parallel implications are either tested or explicitly declared out of scope;
  \item mesh responsibilities remain in VoronoiMeshMaker;
  \item no unnecessary implementation complexity leaks into the user API.
\end{enumerate}

\section{Summary}

The Mohid-NG code standard can be summarised as follows:

\begin{quote}
Use modern C++ internally, expose a simple public API, keep all mesh implementation in VoronoiMeshMaker, operate only on Voronoi-based computational meshes, isolate external backends, design hot kernels around contiguous data, validate public inputs, report errors precisely, test every scientific feature and document every physical assumption.
\end{quote}

This standard is intended to protect the project against three risks: a solver that is fast but unmaintainable, a code base that is flexible but scientifically opaque, and a public API that exposes implementation complexity to users. Mohid-NG shall avoid all three.

\end{document}
